\section{The Requiem Liturgy, Examination of Conscience, and Received Prayers}

The Church does not teach judgment and Purgatory from theological chairs alone; she enacts them whenever she vests in black, lights unbleached candles around a bier, and offers the Holy Sacrifice for the dead. Nowhere are awe before the tribunal and confidence in Christ's mercy bound more closely together than in the Roman Canon and the All Souls' Requiem on 2 November.

In every Mass of the Roman Rite, moments after the Consecration makes Calvary present on the altar, the priest joins his hands in silent recollection of those he commends and prays the \textit{Memento etiam}:

\begin{quote}
\textit{Memento etiam, Domine, famulorum famularumque tuarum N.\ et N., qui nos praecesserunt cum signo fidei, et dormiunt in somno pacis. Ipsis, Domine, et omnibus in Christo quiescentibus, locum refrigerii, lucis et pacis, ut indulgeas, deprecamur.}\endnote{1604 \textit{Missale Romanum}, Canon Missae (\textit{Commemoratio pro defunctis}, unchanged in the 1962 \textit{Missale Romanum}).}
\end{quote}

``Remember also, O Lord, thy servants and handmaids N.\ and N., who have gone before us with the sign of faith, and sleep the sleep of peace. To these, O Lord, and to all that rest in Christ, grant, we beseech thee, a place of refreshment, light, and peace.'' They bear ``the sign of faith'' and ``rest in Christ,'' yet the Church still begs for them \textit{locum refrigerii, lucis et pacis}. Immediately afterward, the priest strikes his breast in the \textit{Nobis quoque peccatoribus}, asking God to admit the living into the saints' company \textit{non aestimator meriti, sed veniae, quaesumus, largitor}---``not weighing our merits, but pardoning our offences.''

On All Souls' Day, the First Mass unfolds this same harmony.\endnote{1962 \textit{Missale Romanum}, In Commemoratione Omnium Fidelium Defunctorum, First Mass (\textit{Ad primam Missam}), Introit, Collect, Gradual, Tract, and Sequence.} Omitting the \textit{Gloria Patri}, the Introit opens with the immemorial cry from 4 Esdras 2:34--35---\textit{Requiem aeternam dona eis, Domine: et lux perpetua luceat eis} (``Eternal rest grant unto them, O Lord, and let perpetual light shine upon them'')---and Psalm 64:2--3 (\textit{ad te omnis caro veniet}, ``all flesh shall come to thee''). The Gradual answers with Psalm 111:7: \textit{In memoria aeterna erit justus: ab auditione mala non timebit} (``The just shall be in everlasting remembrance: he shall not fear the evil hearing'').

Before the Gospel of John 5:25--29 (declaring that the Father gave judgment to the Son of man), the Tract \textit{Absolve, Domine} lifts the Church's plea straight to the tribunal:

\begin{quote}
\textit{Absolve, Domine, animas omnium fidelium defunctorum ab omni vinculo delictorum. V.\ Et gratia tua illis succurrente, mereantur evadere judicium ultionis. V.\ Et lucis aeternae beatitudine perfrui.}\endnote{1962 \textit{Missale Romanum}, Commemoratio omnium fidelium defunctorum, First Mass, Tract.}
\end{quote}

``Absolve, O Lord, the souls of all the faithful departed from every bond of sin. And by the help of thy grace may they be enabled to escape the avenging judgment. And enjoy the bliss of eternal light.''

Then rises the Sequence \textit{Dies irae}. Its first eight strophes set all humanity before the Last Day: the world dissolving into ashes, trembling when the Judge comes to examine all strictly (\textit{Cuncta stricte discussurus}), the trumpet sounding across every tomb, death and nature stunned as creation rises to answer, the written book opened with nothing hidden unavenged---until the seventh strophe faces human poverty (\textit{Cum vix justus sit securus?}) and the eighth casts the soul upon sheer grace: \textit{Rex tremendae majestatis, / Qui salvandos salvas gratis, / Salva me, fons pietatis}.

In the Office of the Dead and at the burial threshold, the Church places Psalm 129---\textit{De profundis}---on our lips.\endnote{\textit{The Little Office of the Blessed Virgin Mary and the Office of the Dead} (New York: Benziger Brothers, 1915), pp.\ 316, 378.} Its antiphon asks starkly, \textit{Si iniquitates observaveris, Domine, Domine, quis sustinebit?} (Ps 129:3), yet the Psalm turns at once to ``merciful forgiveness'' and ``plentiful redemption'' (Ps 129:4, 7--8). As St.\ Robert Bellarmine observes on Psalm 129, true repentance requires both movements together: ``to reflect on our own wretched condition, and to know the extent of God's mercy; because he that is ignorant of the state he is in, seeks for no medicine, does no penance; and he that has no idea of God's mercy, falls into despair, and looks upon penance as of no value.''\endnote{St.\ Robert Bellarmine, \textit{Commentary on the Book of Psalms}, on Ps 129:4--5, trans.\ O'Sullivan (1866).}

\subsection*{Examination of Conscience}

The \textit{Baltimore Catechism No.\ 3} teaches that we prepare daily for judgment ``by a good examination of conscience, in which we will discover our sins and learn to fear the punishment they deserve.''\endnote{\textit{Baltimore Catechism No.\ 3} (1921), Q.\ 1377, p.\ 307.} Let the soul therefore set its own life plainly beneath the light of Christ:

\begin{itemize}
\item \textit{Readiness for the Particular Judgment.} If the Lord summoned me tonight to stand alone before Him, what unconfessed sin, hidden duplicity, or neglected duty of my state would stand exposed? Do I examine my conscience nightly?
\item \textit{Deliberate venial sin and attachments.} What habitual venial faults---sharp speech, gossip, vanity, impatience, slackness in prayer, petty resentment, or attachment to comfort---am I calmly tolerating as ``wood, hay, and stubble''?
\item \textit{Satisfaction and restitution.} After absolution, do I perform my penance reverently, repair harm done to another's property or good name, and embrace prayer, fasting, alms, and daily crosses to heal the wounds of sin?
\item \textit{Works of mercy and lasting example.} Measured by Matthew 25, whom have I fed, welcomed, visited, consoled, or forgiven among ``these my least brethren,'' and what lasting fruit will my words and example leave behind after death?
\item \textit{Charity for the Holy Souls.} Do I assist departed parents, relatives, friends, priests, and forgotten souls in Purgatory through Masses, daily prayer, alms, and indulgences offered by way of suffrage?
\end{itemize}

\subsection*{Received Prayers}

To turn doctrine into worship and intercession, the Church hands down words tested by centuries at the altar and in quiet prayer.

\subsubsection*{Opening Strophes of the Sequence \textit{Dies irae}}

\textbf{Latin (1962 \textit{Missale Romanum}).}

\begin{quote}
\textit{Dies irae, dies illa,\\
Solvet saeclum in favilla:\\
Teste David cum Sibylla.}

\textit{Quantus tremor est futurus,\\
Quando Judex est venturus,\\
Cuncta stricte discussurus!}

\textit{Tuba, mirum spargens sonum\\
Per sepulcra regionum,\\
Coget omnes ante thronum.}

\textit{Mors stupebit, et natura,\\
Cum resurget creatura,\\
Judicanti responsura.}

\textit{Liber scriptus proferetur,\\
In quo totum continetur,\\
Unde mundus judicetur.}

\textit{Judex ergo cum sedebit,\\
Quidquid latet, apparebit:\\
Nil inultum remanebit.}

\textit{Quid sum miser tunc dicturus?\\
Quem patronum rogaturus,\\
Cum vix justus sit securus?}

\textit{Rex tremendae majestatis,\\
Qui salvandos salvas gratis,\\
Salva me, fons pietatis.}\endnote{Latin text from the 1962 \textit{Missale Romanum}, Commemoratio omnium fidelium defunctorum, First Mass, Sequence \textit{Dies irae}, strophes 1--8.}
\end{quote}

\textbf{English (W.\ J.\ Irons).}

\begin{quote}
Day of Wrath! O day of mourning!\\
See fulfill'd the prophets' warning!\\
Heav'n and earth in ashes burning!

Oh, what fear man's bosom rendeth\\
When from Heav'n the Judge descendeth,\\
On Whose sentence all dependeth!

Wondrous sound the trumpet flingeth,\\
Through earth's sepulchres it ringeth,\\
All before the Throne it bringeth.

Death is struck, and nature quaking,\\
All creation is awaking,\\
To its Judge an answer making.

Lo! the Book exactly worded,\\
Wherein all hath been recorded;\\
Thence shall judgment be awarded.

When the Judge His seat attaineth,\\
And each hidden deed arraigneth,\\
Nothing unavenged remaineth.

What shall I, frail man, be pleading,\\
Who for me be interceding,\\
When the just are mercy needing?

King of Majesty tremendous,\\
Who dost free salvation send us,\\
Fount of pity, then befriend us!\endnote{Received English verse translation by W.\ J.\ Irons, as printed in \textit{Hymns Ancient and Modern}, Complete Edition (London: William Clowes and Sons, Limited, catalogued 1890), Hymn 398, p.\ 307, stanzas 1--4.}
\end{quote}

\subsubsection*{Collect for All the Faithful Departed (\textit{Fidelium, Deus})}

\textbf{Latin (1962 \textit{Missale Romanum} and \textit{Office of the Dead}).}

\begin{quote}
\textit{Fidelium, Deus, omnium conditor et redemptor: animabus famulorum famularumque tuarum remissionem cunctorum tribue peccatorum; ut indulgentiam, quam semper optaverunt, piis supplicationibus consequantur: Qui vivis et regnas cum Deo Patre in unitate Spiritus Sancti, Deus, per omnia saecula saeculorum. R.\ Amen.}\\
\textit{V.\ Requiem aeternam dona eis, Domine.}\\
\textit{R.\ Et lux perpetua luceat eis.}\\
\textit{V.\ Requiescant in pace.}\\
\textit{R.\ Amen.}\endnote{Collect body from the 1962 \textit{Missale Romanum}, Commemoratio omnium fidelium defunctorum, First Mass, Collect (corroborated in \textit{A Manual of Prayers for the Use of the Catholic Laity}, New York: Catholic Publication Society, 1889, p.\ 269); full liturgical conclusion and closing versicles from \textit{The Little Office of the Blessed Virgin Mary and the Office of the Dead} (New York: Benziger Brothers, 1915), pp.\ 278--280.}
\end{quote}

\textbf{English (\textit{Office of the Dead}, 1915).}

\begin{quote}
O God, the Creator and Redeemer of all the faithful, give to the souls of Thy servants departed, men and women, the remission of all their sins, that by pious supplications they may attain the pardon which they have always desired. Who livest and reignest, with the Father, in the unity of the Holy Ghost, God, world without end. R.\ Amen.\\
V.\ Grant them eternal rest, O Lord.\\
R.\ And let perpetual light shine on them.\\
V.\ May they rest in peace.\\
R.\ Amen.\endnote{\textit{The Little Office of the Blessed Virgin Mary and the Office of the Dead} (New York: Benziger Brothers, 1915), pp.\ 279--281; cf.\ \textit{A Manual of Prayers for the Use of the Catholic Laity} (1889), pp.\ 269, 775.}
\end{quote}

\subsubsection*{Prayer of St.\ Alphonsus Maria de Liguori on Judgment}

\begin{quote}
Ah my Jesus and my Redeemer, who shall be one day my Judge, pardon me before that day arrives. ``\textit{Non avertas faciem tuam a me.}'' ``Turn not away thy face from me''---Ps.\ xxvi.\ 9. You are now a father to me; and, like a father, receive into your friendship a son who casts himself with sorrow at your feet. My Father, I ask pardon. I have offended you; I have unjustly forsaken you. You did not deserve such treatment from me. I repent of it; I am sorry for it with my whole heart. ``Turn not away thy face from me;'' do not cast me off as I deserved. Remember the blood which you have shed for me, and have pity on me. My Jesus, I wish for no other judge than you. ``I willingly,'' said St.\ Thomas of Villanova, ``submit to the judgment of him who died for me, and who, that he might not condemn me, has condemned himself to the cross.'' St.\ Paul has said the same. ``Who is he that shall condemn? Christ Jesus, that died''---Rom.\ viii.\ 34. My Father, I love you, and, for the future, I will never more depart from your feet. Forget the injuries I have done you, and give me a great love for your goodness. I desire to love you more than I have offended you; but, if you do not assist me, I shall not be able to love you. Assist me, O my Jesus; make me always grateful to your love, that, on the day of judgment, I may be found in the valley of Josaphat, among the number of your lovers. O Mary, my queen and my advocate, assist me now; for if I am lost, you shall not be able to help me on that day. You pray for all; pray also for one who glories in being your devoted servant, and who places so much confidence in you.\endnote{St.\ Alphonsus Maria de Liguori, \textit{Preparation for Death}, Consideration XXV, First Point, ``Affections and Prayers,'' trans.\ by a Catholic clergyman (Dublin: James Duffy, 1843), pp.\ 303--304.}
\end{quote}
